\section{ที่มาและความสำคัญของโครงงาน (Background and Significance)}
ในปัจจุบัน กระแสการดูแลสุขภาพและการออกกำลังกายได้รับความนิยมเพิ่มขึ้นอย่างต่อเนื่อง ส่งผลให้ธุรกิจฟิตเนสหรือยิม (Gym) มีการเติบโตและขยายตัวอย่างรวดเร็ว อย่างไรก็ตาม การบริหารจัดการยิมแบบดั้งเดิมที่อาศัยการจดบันทึกด้วยกระดาษหรือการใช้โปรแกรมตารางคำนวณ (Spreadsheet) ทั่วไป มักประสบปัญหาความยุ่งยาก ซับซ้อน และขาดประสิทธิภาพ โดยเฉพาะอย่างยิ่งในการบริหารจัดการตารางการใช้งานอุปกรณ์ การนัดหมายเทรนเนอร์ และการจัดเก็บข้อมูลสมาชิก ซึ่งอาจนำไปสู่ข้อผิดพลาดต่างๆ เช่น การจองเวลาซ้อนทับกัน (Double Booking), ข้อมูลสมาชิกสูญหาย หรือความล่าช้าในการตรวจสอบสถานะของอุปกรณ์และเทรนเนอร์

เพื่อแก้ไขปัญหาดังกล่าวและยกระดับการให้บริการ คณะผู้จัดทำจึงได้แนวคิดในการพัฒนา ``ระบบบริหารจัดการยิม (GymBro Management System)'' ซึ่งเป็นเว็บแอปพลิเคชันที่ถูกออกแบบมาเพื่อช่วยอำนวยความสะดวกในการดำเนินงานของยิมอย่างครบวงจร ระบบนี้จะทำหน้าที่เป็นศูนย์กลางในการเชื่อมโยงข้อมูลระหว่างสมาชิก (Members), ผู้ดูแลระบบ (Admins), เทรนเนอร์ (Trainers), และทรัพยากรภายในยิม (Equipments) เข้าด้วยกัน ช่วยให้การบริหารจัดการเป็นไปอย่างเป็นระบบ รวดเร็ว และตรวจสอบได้ ซึ่งไม่เพียงแต่จะช่วยลดภาระงานของผู้ดูแลระบบ แต่ยังช่วยสร้างประสบการณ์ที่ดีให้กับสมาชิกผู้ใช้บริการอีกด้วย

\section{วัตถุประสงค์ของโครงงาน (Objectives)}
\begin{enumerate}
    \item เพื่อพัฒนาเว็บแอปพลิเคชันสำหรับการบริหารจัดการข้อมูลและการจองคลาสเรียนในฟิตเนส
    \item เพื่ออำนวยความสะดวกให้กับสมาชิกในการตรวจสอบตารางเวลาและจองคลาสเรียนด้วยตนเอง
    \item เพื่อช่วยให้ผู้ดูแลระบบสามารถบริหารจัดการข้อมูลสมาชิก เทรนเนอร์ และอุปกรณ์ได้อย่างมีประสิทธิภาพ
    \item เพื่อป้องกันปัญหาการจองเวลาซ้อนทับกัน (Double Booking) ของเทรนเนอร์และอุปกรณ์
    \item เพื่อศึกษาและประยุกต์ใช้ความรู้ด้าน Full Stack Development ในการพัฒนาระบบใช้งานจริง
\end{enumerate}

\section{ขอบเขตของโครงงาน (Scope)}
โครงงานนี้ครอบคลุมการพัฒนาระบบเว็บแอปพลิเคชัน โดยแบ่งขอบเขตการทำงานออกเป็น 2 ส่วนหลัก ดังนี้:

\subsection{ส่วนของผู้ดูแลระบบ (Administrator)}
\begin{itemize}
    \item สามารถดูภาพรวมสถิติของระบบ (Dashboard) เช่น จำนวนสมาชิก, เทรนเนอร์, อุปกรณ์ และรายการจองในปัจจุบัน
    \item สามารถจัดการข้อมูลสมาชิก (เพิ่ม, ลบ, แก้ไข, ตรวจสอบข้อมูล)
    \item สามารถจัดการข้อมูลเทรนเนอร์ (เพิ่ม, ลบ, แก้ไข, ระดับความสามารถ, ความเชี่ยวชาญ)
    \item สามารถจัดการข้อมูลอุปกรณ์ออกกำลังกาย (เพิ่ม, ลบ, แก้ไขสถานะการใช้งาน/ซ่อมบำรุง)
    \item สามารถบริหารจัดการตารางการฝึกซ้อม (Session Management) และตรวจสอบประวัติการจอง
    \item สามารถดูบันทึกการใช้งานระบบ (System Logs) เพื่อตรวจสอบการกระทำย้อนหลังได้
\end{itemize}

\subsection{ส่วนของสมาชิก (Member)}
\begin{itemize}
    \item สามารถลงทะเบียนสมัครสมาชิกและเข้าสู่ระบบได้
    \item สามารถดูตารางเวลาการฝึกซ้อม (Schedule) ของยิมได้
    \item สามารถจองเวลาการฝึกซ้อม (Booking) โดยเลือกเทรนเนอร์และอุปกรณ์ที่ต้องการได้
    \item ระบบจะตรวจสอบความว่างของเทรนเนอร์และอุปกรณ์ให้โดยอัตโนมัติก่อนการจอง
    \item สามารถยกเลิกการจองของตนเองได้
    \item สามารถดูและตรวจสอบข้อมูลโปรไฟล์ของตนเองได้
\end{itemize}

\section{เป้าหมายของโครงงาน (Goals)}
\begin{enumerate}
    \item ได้ระบบบริหารจัดการยิมที่สามารถใช้งานได้จริง มีความเสถียร และถูกต้องตามความต้องการ
    \item ระบบมีความปลอดภัยในระดับพื้นฐาน มีการตรวจสอบสิทธิ์การเข้าใช้งาน (Authentication & Authorization)
    \item ส่วนติดต่อผู้ใช้งาน (User Interface) มีความทันสมัย ใช้งานง่าย (User-friendly) และรองรับการแสดงผลบนอุปกรณ์ต่างๆ (Responsive)
    \item โค้ดของโปรแกรมมีความเป็นระเบียบ ง่ายต่อการดูแลรักษาและพัฒนาต่อยอด
\end{enumerate}

\section{ประโยชน์ที่คาดว่าจะได้รับ (Expected Benefits)}
\begin{enumerate}
    \item ช่วยลดความผิดพลาดในการบริหารจัดการตารางเวลาและการจองคลาสเรียน
    \item เพิ่มความสะดวกสบายให้กับสมาชิกในการเข้าถึงบริการของยิม
    \item ผู้ดูแลระบบสามารถทำงานได้รวดเร็วและมีประสิทธิภาพมากขึ้น ลดการใช้กระดาษและขั้นตอนที่ซ้ำซ้อน
    \item มีฐานข้อมูลที่เป็นระบบ สามารถนำข้อมูลไปวิเคราะห์เพื่อวางแผนการบริหารงานในอนาคตได้
    \item ผู้จัดทำได้รับความรู้และประสบการณ์ในการพัฒนา Full Stack Web Application ตั้งแต่การออกแบบจนถึงการ Deploy
\end{enumerate}
